\documentclass[11pt]{article}
\usepackage[margin=25mm]{geometry}
\usepackage{fontspec}
\setmainfont{Latin Modern Roman}
\usepackage{amsmath,amssymb,mathtools}
\usepackage{unicode-math}
\setmathfont{Latin Modern Math}
\usepackage{xurl,hyperref}
\hypersetup{colorlinks=true,urlcolor=blue}
\setcounter{secnumdepth}{0}
\setlength{\parindent}{0pt}
\setlength{\parskip}{5pt}
\setlength{\emergencystretch}{3em}
\newcommand{\sref}[2]{\hyperlink{source:#1:#2}{[S#2]}}
\newcommand{\cataloguescope}[1]{\paragraph{Recorded target.}#1}
\newcommand{\R}{\mathbb R}
\newcommand{\C}{\mathbb C}
\newcommand{\N}{\mathbb N}
\newcommand{\Z}{\mathbb Z}
\newcommand{\Q}{\mathbb Q}
\newcommand{\F}{\mathbb F}
\newcommand{\eps}{\varepsilon}
\begin{document}
\section{Equality of Capacity–Volume and ADM Mass}\label{statement}

\subsection{Definitions and mathematical statement}

\paragraph{Shared notation.}$\mathbb N=\{1,2,\ldots\}$; the standard number fields have their usual meanings. All additional conventions are stated locally.


An asymptotically flat three-manifold approaches Euclidean space on its end. The ADM mass is the asymptotic flux of first derivatives of the metric. The normalized capacity of a compact region $K$ is the infimum of $(4\pi)^{-1}\int_{M\setminus K}|\nabla u|^2dV_g$ over admissible functions with boundary value one on $K$ and value zero at infinity. Writing $|K|$ for Riemannian volume, Jauregui's capacity--volume mass is
\[
m_{\mathrm{CV}}(M,g)=\sup_{\{K_j\}}\limsup_{j\to\infty}
\left[\left(\frac{3|K_j|}{4\pi}\right)^{1/3}-\operatorname{cap}(K_j)\right],
\]
where the supremum is over the admissible compact exhaustions in the source.

\paragraph{Catalogue wording retained for traceability.}
For every asymptotically flat Riemannian three-manifold $(M^3,g)$, does the capacity–volume mass equal the ADM mass, $m_{\mathrm{CV}}(M,g)=m_{\mathrm{ADM}}(M,g)$?

\paragraph{Scope and source qualification.}The source reports equality for harmonically flat ends under additional hypotheses, and asks about general asymptotically flat manifolds. An exact proof must use the report's definition of the exhaustion supremum, rather than a single chosen family of regions.
\subsection{Short English statement}
Determine whether the mass defined from large-region capacity and volume always equals ADM mass.
\subsection{Research attempt}
\paragraph{Model, date and investigated scope.}
OpenAI GPT-6 Astra (configured model \texttt{gpt-6-astra}), ultra reasoning; generated 11 October 2026. Computed the round-exhaustion capacity–volume limit and ADM mass on a harmonic conformal end; did not control the exhaustion supremum or general ends.

\paragraph{Sources and literature check.}
Read the supplied definition, catalogue wording, source qualifications and reference entries for this question. The historical source/context and inherited status are retained; no fresh literature or novelty certification is claimed. No external theorem is used without proof in this attempt; explicit model assumptions are distinguished from facts about the source model.

\paragraph{A harmonic-end exhaustion calculation.} Let \(g=(1+a/r)^4\delta\) outside a compact set in three dimensions, with \(a\ge0\). For the coordinate sphere \(r=R\), a radial capacitary potential \(v\) satisfies \(((r+a)^2v')'=0\), \(v(R)=1\), \(v(\infty)=0\). Thus \(v(r)=(R+a)/(r+a)\). Its normalized energy is
\[
 \operatorname{cap}(B_R)=\int_R^\infty(r+a)^2|v'|^2dr=R+a.
\]
It minimizes the energy: for any zero-boundary perturbation \(w\), the cross term vanishes by the radial harmonic equation and integration by parts, leaving the nonnegative energy of \(w\).
Meanwhile the compact interior contributes only a bounded volume correction, and expansion of \(r^2(1+a/r)^6\) gives
\[
 |B_R|=4\pi(R^3/3+3aR^2+O(R)),\qquad
 (3|B_R|/(4\pi))^{1/3}=R+3a+O(R^{-1}).
\]
The capacity--volume difference along this exhaustion tends to \(2a\). The ADM flux also equals \(2a\): since \(g_{ij}=(1+4a/r+O(r^{-2}))\delta_{ij}\), its flux integrand is \(-2\partial_r(1+a/r)^4\), and division by \(16\pi\) gives the stated limit.
\paragraph{The supremum gap.} This reconstructs the expected value along round regions in a known special end. The definition takes a supremum over all admissible exhaustions; one chosen exhaustion only gives a lower bound for that supremum. Bounding arbitrary distorted exhaustions, and controlling non-harmonic asymptotically flat errors, are the first missing steps. The source already reports a harmonic-end equality theorem under additional assumptions, so this calculation is not new progress toward the general equality.

\paragraph{Verification and adversarial audit.}
Checked the harmonic potential, energy minimization, volume coefficient and ADM flux; kept a single exhaustion distinct from the source supremum. This is a GPT-6 Astra Ultra self-audit, not independent human verification. Prompt A drafting followed by prompt B scope/dependency/boundary self-audit. Checked the harmonic potential, energy minimization, volume coefficient and ADM flux; kept a single exhaustion distinct from the source supremum. The remaining gap is: Bound every admissible exhaustion and control general asymptotically flat errors beyond harmonic conformal ends. Unresolved status and zero original-target completion are retained.

Coordinator checked the radial capacitary potential, normalized energy, volume expansion and ADM coefficient; the supremum over all exhaustions is left open. This additional check was model-assisted, not independent human verification.

\paragraph{Tools, process and contributions.}
Python 3 and local file-reading/writing tools for manuscript serialization and structural checks; mathematical calculations were symbolic reasoning, with no numerical experiment or theorem prover. Contribution prompt A was followed by prompt B. GPT-6 Astra with ultra reasoning drafted the research and performed its stated self-audit; any additional coordinator contribution is identified in the verification and audit. A GPT-6-based Codex coordinator assembled the manuscript and checked repository integration; its exact variant and reasoning configuration were not exposed. No human mathematical contribution, independent human proof check or formal proof certification is supplied. Complete derivations and the consulted source references are included here; private reasoning transcripts are not published.

\paragraph{Final status and first remaining gap.}
\textbf{UNRESOLVED}. Bound every admissible exhaustion and control general asymptotically flat errors beyond harmonic conformal ends.

\section{Completion Estimate}
\noindent\textbf{COMPLETION: 0\%}
The proved calculation is an elementary auxiliary result or a test of a restricted strategy; no new portion of the full source target is established.

\subsection{Sources}
\begin{itemize}
\item[\hypertarget{source:statement:1}{S1}] ULAM.AI and the UnsolvedMath Contributors, UnsolvedMath, catalogue entry. Catalogue curation under CC BY 4.0; definitions and English exposition added locally.
\url{https://huggingface.co/datasets/ulamai/UnsolvedMath}.
\item[\hypertarget{source:statement:2}{S2}] Mathematical Aspects of General Relativity (2021), Oberwolfach Reports, DOI 10.4171/owr/2021/40, pp. 2157--2267. Original report and mathematical context.
\url{https://doi.org/10.4171/owr/2021/40}.
\end{itemize}
\label{end:statement}
\end{document}
